%============ ORIENTATION ============
\section{Orientation: what Part VI is actually for}
\markboth{Orientation}{Orientation}

Parts I--V of \textit{Where Mathematics Comes From} build an apparatus. \textbf{Part VI is the stress
test.} Lakoff \& Núñez take the single most revered formula in classical mathematics and ask whether
their method --- \emph{mathematical idea analysis} --- can say something about it that axioms,
definitions, theorems and proofs cannot.

The opening move is rhetorical but surgical: Benjamin Peirce, one of Harvard's leading nineteenth\-/century
mathematicians, lecturing on Euler's proof and confessing that he had proved something he did not
understand. Part VI is a sustained argument that \textbf{proof is not understanding}, and that
understanding has a describable structure built from three mechanisms the book has already
introduced.

\begin{center}
\begin{tikzpicture}[font=\footnotesize,
   mech/.style={rounded corners=4pt,draw=#1,line width=1pt,fill=#1!7,text width=4.55cm,
                align=left,inner sep=8pt,minimum height=2.35cm,anchor=north west}]
  \node[mech=purple] (a) at (0,0)
    {\textbf{\color{purple}Grounding metaphors}\\[3pt]
     Map \emph{bodily, pre-mathematical} experience onto mathematics.\\[3pt]
     \textit{e.g.} \meta{Arithmetic Is Object Construction}};
  \node[mech=skyblue] (b) at (5.05,0)
    {\textbf{\color{skyblue}Linking metaphors}\\[3pt]
     Map one \emph{mathematical} domain onto another, exporting inference.\\[3pt]
     \textit{e.g.} \meta{Functions Are Numbers}};
  \node[mech=teal] (c) at (10.1,0)
    {\textbf{\color{teal}Conceptual blends}\\[3pt]
     Fuse two domains into a third with \emph{emergent} entities and entailments.\\[3pt]
     \textit{e.g.} the \blend{Unit Circle blend}};
  \node[anchor=north,text width=14.6cm,align=center,font=\scriptsize\itshape,text=slate]
     at (7.325,-2.62) {plus the \textbf{BMI}, the Basic Metaphor of Infinity, wherever a limit appears};
\end{tikzpicture}
\end{center}

\vspace{2pt}
\subsection{The standard of success they set themselves}

This is worth copying into the front of your notebook, because it is the yardstick against which the
rest of the Part must be measured.

\begin{keybox}[title={Requirements for an adequate mathematical idea analysis \normalfont\itshape(book p.~436 / pdf p.~455)}]
Of an equation, definition, axiom or theorem, an idea analysis must give:
\begin{enumerate}[label=\textbf{\arabic*.}]
  \item the \textbf{metaphor-and-blend structure} that makes it true \emph{relative to that structure};
  \item the \textbf{conceptual relationships} among the elements that appear in, or are presupposed by, it;
  \item the \textbf{ideas expressed by the symbols};
  \item the \textbf{grounding} of the concepts used in the analysis.
\end{enumerate}
\end{keybox}

\noindent The consequence they draw immediately is the most philosophically loaded sentence in Part VI:

\begin{bookquote}{book p.~436 / pdf p.~455}
\textit{The ``truth'' of an equation, axiom, or theorem depends on what it means. That is, truth is
relative to a conceptual structure\ldots\ Thus, $e^{\pi i}+1=0$ is not true of Euclidean geometry or
the arithmetic of real numbers for an obvious reason: Those mathematical domains do not have the
right idea structure to give meaning to the equation, much less to make it true.}
\end{bookquote}

\noindent I return to this claim --- and to the two very different readings it admits --- in
\S\ref{sec:audit}.

\subsection{Why this particular equation}

Their own justification, stated on the Part-opening page, is that Euler's equation is a
\emph{junction}: it is the one place where every branch of classical mathematics is forced into
contact.

\begin{bookquote}{book p.~381 / pdf p.~400}
\textit{It brings together in one equation the most important numbers in mathematics: $e$, $\pi$,
$i$, $1$, and $0$. In doing so, this one equation implicitly states a relationship among the
branches of classical mathematics---arithmetic, algebra, geometry, analytic geometry, trigonometry,
calculus, and complex variables.}
\end{bookquote}

\noindent And their statement of the problem is refreshingly blunt: the usual gloss on exponentiation
simply \emph{cannot} be what is going on.

\begin{warmbox}[title={The problem, stated precisely \normalfont\itshape(book p.~384 / pdf p.~403)}]
$e^{\pi i}$ makes no sense as ``$e$ multiplied by itself $\pi$ times, and then $i$ times.''
$\pi$ is an infinite decimal; $i=\sqrt{-1}$. \textbf{``What could `$\sqrt{-1}$ times' possibly mean?''}
A different answer is required --- one that follows from the \emph{conceptual content} of
$e$, $\pi$, $i$ and $-1$.
\end{warmbox}
